\question
Первокурсники ИСа отмечают темы, которые будут на консультации по дму, из общего списка. $U = \{1, 2, 3, 4, 5, 6, 7\}$ -- список всех тем. $A_1$, $A_2$, $A_3$, $A_4$, $A_5$, $A_6$ -- множества тем, отмеченных студентами 1, 2, 3, 4, 5, 6 соответственно.
\\
\begin{align*}
    A_1 &= \{1, 3\}, & A_2 &= \{3, 4\}, & A_3 &= \{2, 4\}, \\
    A_4 &= \{1, 3, 4, 5, 6\}, & A_5 &= \{1, 3\}, & A_6 &= \{5, 7\}.
\end{align*}
На множестве студентов $\{1, 2, 3, 4, 5, 6\}$ определены два отношения:
\begin{enumerate}
    \renewcommand{\labelenumi}{\arabic{enumi})}
    \item $iRj$, если $A_i \subseteq A_j$ (темы студента $i$ полностью содержатся в темах студента $j$);
    \item $iQj$ (при $i \ne j$), если $A_i \cap A_j \ne \varnothing$ (у студентов $i$ и $j$ есть хотя бы одна общая тема).
\end{enumerate}
Вам, как ментору по дму, чтобы понять, что нужно покрыть на консультации, нужно сделать следующее:
\begin{enumerate}
    \renewcommand{\labelenumi}{\alph{enumi})}
    \item Определите свойства $R$ (рефлексивность, транзитивность, симметричность, антисимметричность). Для каждого свойства приведите доказательство или контрпример.
    \item Определите свойства $Q$ (рефлексивность, транзитивность, симметричность, антисимметричность). Для каждого свойства приведите доказательство или контрпример.
    \item Найдите единственную пару $(p, 6)$, в которой участвует студент 6, для которой $pQ6$ или $6Qp$ выполняется, но при этом между ними не существует ни $pR6$, ни $6Rp$.
\end{enumerate}

---------------

Автор -- Денис Алдатов
